\section{Discussion}

This paper contributes the most recent estimate of prepregnancy underweight for India and the first estimates stratified by social group. We find that in 2019-2021, about one in five women in India was underweight at the beginning of pregnancy. This represents significant progress; the prevalence of prepregnancy underweight has come down by over 50\% since 2005.  However, this indicator suggests that maternal nutrition remains an important population health problem.\footnote{For comparison, in the United States, fewer than 5\% of women are underweight at the start of pregnancy \citep{zong2022maternal}.} 

Within India, the burden of maternal undernutrition falls disproportionately on Adivasi, Dalit, and OBC women, who have a prevalence of prepregnancy underweight 11, 7, and 5 percentage points higher than forward caste Hindus, respectively.  The prevalence of prepregnancy underweight among Forward caste Hindu and Muslim women is similar.  This result coheres with other research that finds advantages in Muslim women's and children's health relative to Hindus' despite social exclusion.   

The Kitagawa decompositions show that higher fertility and more closely spaced births among Adivasis, Dalits, and OBCs do not explain why they are more likely to be underweight before pregnancy than Forward caste Hindu women.  This is in part because the differences by social group in parity and birth spacing are small, and in part because these variables are only weakly correlated with prepregnancy underweight in the most recent Indian data. Differences in household wealth, in contrast, do explain large proportions of the gaps.  Future studies might investigate whether this is because wealthier households are able to afford food, because wealthier households can afford labor saving devices such as vehicles or washing machines, or because wealthier households offer better protection against disease.  It is also possible that wealthier households have more equitable intrahousehold food allocation.

We note that a significant fraction of the gaps in the prevalence of prepregnancy underweight remains unexplained after accounting for differences in household wealth.  Could remaining gaps in the prevalance of prepregnancy underweight be explained by differences in intrashousehold allocation of food across caste groups? Unfortunately, the NFHS does not have information about the relative quantities of food consumed by different household members.  However, the literature and some evidence from the IDHS suggest that this may not be important for explaining the remaining gaps in the prevalence of prepregnany underweight.  \cite{harris2017determinants}'s review  of intrahousehold food allocation among adults in India suggests that oppressed caste households may distribute food more equally between men and women than Forward caste Hindu households, perhaps because oppressed caste women are more likely to work outside the home than dominant caste women \citep{chakrabarty1996gender, eswaran2013status}.  To provide some preliminary evidence on this question, we analyzed an IHDS (2011) question which asked ever-married women whether women eat with, or after, the men in their households.  Figure \ref{fig:ihds_eat_last} in the online appendix shows the proportion of women who report eating last, is very similar across social group.\footnote{\cite{desai2014muslim} analyzed IHDS (2005) data and found that Muslim women \emph{were} less likely to report eating after the men in their household than Hindu women.}  

We note that the literature on racial differences in birth outcomes in the United States increasingly points to racism-related stress among mothers as an important reason for higher rates of preterm birth and lower birth weights among Black infants \citep{hogue2005stress, braveman2015role, sosnaud2021cross}. Could differences in discrimination-related stress explain the remaining caste gaps in maternal undernutrition? Unfortunately, to our knowledge, no population-level study of pregnant women contains data on the biomarkers such as cortisol and C-reactive protein that would help answer that question. 

Our research has several implications for policy and research.  The finding that maternal undernutrition remains high at the population level underscores the importance of the 2013 National Food Security Act's maternity entitlements to pregnant women. The law promised nearly all pregnant women a Rs. 6,000 benefit, but in practice, only a fraction of the funds that would be needed for a universal program have been allocated.\footnote{In 2020, allocation to maternity entitlements was only enough to cover about 1 in 7 women who gave birth \cite{dreze2021derailed}.  Maternity entitlements are often restricted to the first child and are conditional on the woman receiving care from government health clinics \citep{kalra2020birth}.}  If fully funded, indexed for inflation, and given at the start of pregnancy, cash transfers during pregnancy could help address India's remaining burden of prepregnancy underweight \citep{coffey2016underweight}. 

Beyond the National Food Security Act, programs and policies to address maternal undernutrition should also focus on adolescents and never-married women.  The time  between marriage and conception is short: on average, women who had their first births in the three years before the NFHS-5 were married for only 17.23 months before conception.  Cash or in-kind transfers to adolescent girls and young women in their natal households may improve prepregnancy health.  

Finally, this study highlights the value of representative, longitudinal monitoring data on health in pregnancy for India.  The cross sectional data that we analyze permit the construction of synthetic cohorts that yield valuable insights into maternal nutrition, and social group differences in maternal nutrition.  Yet, they do not allow for individual level analyses.  For example, we do not know what fraction of women who begin pregnancy underweight compensate for low prepregnancy body mass with adequate weight gain.  Nor do we know how many women who begin pregnancy underweight go on to have a preterm or low birthweight infant.  Investments in longitudinal health data will improve the health of both mothers and children.

\section{OLD DISCUSSION}

This paper contributes the most recent estimates of the prevalence of prepregnancy underweight for India and the first estimates stratified by social group. Notably, we find that in 2019-2021, about one in five women in India was underweight at the beginning of pregnancy. This represents progress in maternal health; the prevalence of prepregnancy underweight has come down by over 50\% since 2005.  However, the prevalence of prepregnancy underweight remains high.  For comparison, in the United States, fewer than 5\% of women are underweight at the start of pregnancy \citep{zong2022maternal}. Even among socially dominant forward caste Hindus, the burden of prepregnancy underweight remains similar to what it was among all sub-Saharan African women two decades ago \citep{coffey2015prepregnancy}. Previous work on prepregnancy underweight in India has based estimates on the sample of all nonpregnant women of reproductive age, without adjusting for selection into pregnancy -  \cite{miller2025preconception} estimates a prevalence of preconception malnutrition of 14\% for India in 2019-21; our work demonstrates that this is an underestimate. 

Within India, the burden of maternal undernutrition falls disproportionately on Adivasi, Dalit, and OBC women, who have a prevalence of prepregnancy underweight 11, 7, and 4 percentage points higher than forward caste Hindus, respectively. For reference, we find that the gap between the poorest and richest wealth quartile is 18 percentage points. It is also insightful to compare these results to social group inequalities in other morbidities. In the same dataset, disparities in the prevalence of stunting in children and anemia among married women who recently gave birth showed disparities across social groups of similar magnitude to those we find in prepregnancy underweight in this paper \citep{rao2023trends, das2024prevalence}. On the other hand, pregnancy risk classifications—whether measured during or after pregnancy—exhibit much smaller caste stratification than prepregnancy nutritional status \citep{kuppusamy2023high, tandon2023adverse}. We are not aware of prior work that compares prepregnancy underweight by social group in other countries to benchmark these disparities internationally.

While these patterns are consistent with broader social and economic disadvantage, differences in prevalence are likely not due to worse intrahousehold food allocation for women in marginalized caste groups. \cite{harris2017determinants}'s review article of intrahousehold food allocation among adults suggests that that oppressed caste households that are ``not severely food insecure'' may distribute food more equally between men and women than Forward caste Hindu households. This may be because oppressed caste women are more likely to work outside the home than dominant caste women, and so have greater autonomy within their homes \citep{chakrabarty1996gender, eswaran2013status}.  The frequency and nature of women's religious fasting may also differ across social groups.  However, our analysis of an IHDS (2011) question which asked ever-married women whether women eat with, or after, the men in their households, found no difference between Forward Caste Hindus, on the one hand, and Dalits or Adivasis, on the other.  Figure \ref{fig:ihds_eat_last} in the online appendix shows the proportion of women who report eating last, by social group.  Further research is needed to better understand whether there are social group differences in intrahousehold food allocation, and if so, why.  

The Kitagawa decompositions show that higher fertility and more closely spaced births among Adivasis, Dalits, and OBCs do not explain why they are more likely to be underweight before pregnancy than Forward caste Hindu women.  This is in part because the differences by social group in parity and birth spacing are small, and in part because these variables are only weakly correlated with prepregnancy underweight in the most recent Indian data. Differences in household wealth, in contrast, do explain large proportions of the gaps.  Future studies might investigate whether this is because wealthier households are able to afford food, because wealthier households can afford labor saving devices such as vehicles or washing machines, or because wealthier households offer better protection against disease.  

We might also ask whether wealthier households have more gender equitable food ways, independent of caste.  Data from the 1980s suggest that women reduce on food consumption by more than men during times of scarcity to ensure that men have the energy to work \citep{behrman1990intrahousehold}.  It is reasonable to assume that periods of scarcity are more common in poorer households than richer ones.  If so, discriminatory food practices could be more common in poor households than richer ones.  However, there has been substantial economic growth since this study was conducted, and a substantial reduction in the proportion of households reporting food insecurity \citep{deaton2009food}. 

We note that a significant fraction of the gaps in the prevalence of prepregnancy underweight remains unexplained after accounting for differences in household wealth.  We expect that if we were able to measure differences in discrimination-related stress (such as with cortisol) or inflammation (such as with C-reactive protein), that might explain at least part of the remaining gaps.  This is an important hypothesis to consider, as the literature on racial differences in birth outcomes in the United States increasingly points to racism-related stress as an important reason for higher rates of preterm birth and lower birth weights among Black infants \citep{hogue2005stress, braveman2015role, sosnaud2021cross}. Future research should aim to unpack how, and at what cost, childbearing-age women in India experience discrimination based on social group membership.  

Although not the focus of this paper, we note that Adivasi, Dalit, and OBC women do not compensate for prepregnancy underweight with higher weight gain during pregnancy.  The US Institute of Medicine (IOM) recommends that women who begin pregnancy underweight gain 28 to 40 pounds during pregnancy \citep{iom2009weightgain}. Women who have normal or overweight BMIs at the beginning of pregnancy should gain less.\footnote{We cite the IOM guidelines because the Government of India has not issued guidelines on weight gain in pregnancy. The IOM guidelines have also been endorsed by the United Kingdom's National Institute for Health and Care Excellence \citep{NICE2025}.}  Table \ref{tab:addnut} in the online appendix uses NFHS-5 data to show estimates of weight gain during pregnancy for each social group. The methods for deriving these estimates were developed by \cite{coffey2015prepregnancy} and are detailed in section \ref{sec:weightgain} of the online appendix.\footnote{Table \ref{tab:addnut} also computes the prevalence of prepregnancy obesity ($BMI>30$) in India and by social group, using the same method we used to compute the prevalence of prepregnancy underweight.  There are small differences by social group: 11\% of Forward caste Hindu and Muslim women are obese before pregnancy, compared to 5\% of Adivasi women, 7\% of Dalit women, and 8\% of OBC women.  We include these estimates for completeness and to encourage future research on maternal nutrition.  Although obesity rates are increasing in India \citep{luhar2020forecasting}, obesity among women who are likely to become pregnant is relatively uncommon: 8\% of Indian women are estimated to be obese before pregnancy.} There is no statistically significant difference in weight gain during pregnancy across social groups.  For no group does the average estimated weight gain exceed 20 pounds. 

We note the similarity in the prevalence of prepregnancy underweight between Forward caste Hindu and Muslim women.  Despite social and economic disadvantage, Muslim women are no more likely to be underweight before pregnancy, and are less likely to be underweight after accounting for household wealth.  This result coheres with other research that finds advantages in Muslim women's and children's health relative to Hindus'. Could this be because Muslim households allocate food more equitably between young women and other family members than Hindus do?  \cite{desai2014muslim} analyzed IHDS (2005) data and found that Muslim women \emph{were} less likely to report eating after the men in their household than Hindu women.  Figure \ref{fig:ihds_eat_last} in the online appendix shows the same comparison using the IHDS 2011.  In the 2011 data, the difference in the proportion of Muslim and Hindu married women who report eating last was longer statistically significant.  Future qualitative research could find other ways of investigating the relationships among age, gender, and food allocation.

Can policies reduce maternal undernutrition and address inequalities in prepregnancy underweight?  It may be more effective to focus on never-married women than married women.  The time interval between marriage and conception is short: on average, women who had their first births in the three years before the NFHS-5 were married for only 17.23 months before conception.  \cite{banerjee2024persistent} study a cash transfer program in West Bengal in which girls teenage girls received Rs. 1,000 for every year they remained unmarried and in school. They received Rs. 25,000 if they remained in school and unmarried until age 18.  The program delayed marriage and made girls less likely to say that (hypothetical) domestic violence was justified.  Such program might impact girls' nutrition as well.  Because the marginal impact of a fixed amount of cash is greater for poorer girls, such a program would likely improve equity as well.  

Our research also underscores the potential of the 2013 National Food Security Act's maternity entitlements to improve maternal nutrition.  The law promised nearly all pregnant women a Rs. 6,000 benefit, but in practice, only a fraction of the funds that would be needed for a universal program have been allocated.  In 2020, allocation to maternity entitlements was only enough to cover about 1 in 7 women who gave birth \cite{dreze2021derailed}.  Maternity entitlements are often restricted to the first child and are conditional on the woman receiving care from government health clinics \citep{kalra2020birth}.  If fully funded, indexed for inflation, and given at the start of pregnancy, cash transfers during pregnancy could help address India's high burden of prepregnancy underweight \citep{coffey2016underweight}. Policymakers interested in addressing social group differences in maternal nutrition might also consider increasing the amount of the transfer to women from marginalized groups.

Finally, this study highlights the value of representative, longitudinal monitoring data on health in pregnancy for India.  The cross sectional data that we analyze permit the construction of synthetic cohorts that yield valuable insights into maternal nutrition, and social group differences in maternal nutrition.  Yet, they do not allow for individual level analyses.  For example, we do not know what fraction of women who begin pregnancy underweight compensate for low prepregnancy body mass with adequate weight gain.  We do not know how many women who begin pregnancy underweight also suffer from anemia.  We do not know how many women who begin pregnancy underweight have a preterm or low birthweight infant.  Considering that India is now the country with the highest number of births globally, investments in Indian health data will improve global health.